\chapter{АНАЛИЗ ПРЕДМЕТНОЙ ОБЛАСТИ}


\section{Описание предметной области}

\textbf{Класс рассматриваемых задач.} В работе рассматривается класс задач распределения потока пользователей между ресурсами ограниченной вместимости при условии, что инструмент управления распределением — рекомендательная подсказка, не носящая принудительного характера. Цель управляющей стороны — максимизировать суммарную полезность распределения, не допуская при этом ни перегрузки отдельных ресурсов, ни их хронической недозагрузки.

Подобный класс задач возникает в широком спектре прикладных областей. К нему относятся задачи маршрутизации транспортных потоков, в которых навигационный сервис рекомендует водителям маршруты при ограниченной пропускной способности дорожной сети; задачи распределения мест в сервисах бронирования и заведениях обслуживания; задачи управления учебной нагрузкой при ограничениях на численность учебных групп; задачи управления потоками посетителей в зонах с ограниченной пропускной способностью (выставочные пространства, парки, конференции); задачи управления нагрузкой в распределённых вычислительных системах. Несмотря на различия в предметной специфике, перечисленные задачи объединены общей структурой: наличие неоднородного спроса со стороны автономных пользователей, ограниченность ресурсов, отсутствие возможности принудительного назначения и наличие управляющего рычага в виде рекомендательной выдачи.

\textbf{Действующие стороны.} В постановке задачи выделяются три стороны.

Пользователи — субъекты, самостоятельно принимающие решение о выборе ресурса. Множество пользователей характеризуется неоднородностью предпочтений и автономностью поведения.

Ресурсы — объекты выбора. Каждый ресурс характеризуется ограничением по числу пользователей, превышение которого приводит к деградации качества обслуживания: отказам в доступе, ухудшению опыта, снижению эффективности.

Организатор — управляющая сторона, заинтересованная в качестве функционирования системы как целого. Не обладает полномочиями принудительного назначения; влияет на распределение исключительно через формирование рекомендательных подсказок — ранжированных списков, предъявляемых пользователю в качестве предпочтительных вариантов.

\textbf{Особенности задачи.} Управляющий механизм носит вероятностный характер: пользователь не обязан следовать рекомендации. Это отличает задачу от классических задач оптимального назначения, в которых решение управляющей стороны эквивалентно фактическому распределению.

Существенной особенностью является взаимосвязанность рекомендаций. Поскольку пользователи получают подсказки одновременно, рекомендация одному пользователю не может рассматриваться изолированно от рекомендаций, выдаваемых другим. При независимом формировании рекомендаций по индивидуальной релевантности множество пользователей получит в верхних позициях одни и те же объективно наиболее популярные ресурсы, что приведёт к ожидаемому совокупному потоку, превышающему вместимость. Следовательно, оптимальная индивидуальная подсказка для пользователя есть функция от распределения подсказок, выдаваемых остальным пользователям. Задача переходит в плоскость совместной оптимизации.

\textbf{Полезность и противоречие интересов.} Каждое сопоставление пользователя и ресурса характеризуется числовой оценкой полезности — степени соответствия ресурса предпочтениям пользователя с учётом контекстных факторов (времени, заполненности, удобства). Целевая функция организатора — суммарная полезность по всем парам «пользователь–ресурс», подлежащая максимизации при соблюдении ограничений вместимости.

Особенностью класса задач является системное противоречие между индивидуальной и совокупной оптимальностью. Индивидуальный выбор пользователей в пользу наиболее предпочтительного ресурса при отсутствии координации приводит к коллективному результату, заведомо уступающему оптимальному. Управляющая сторона решает задачу нивелирования данного противоречия: смещения индивидуального поведения в направлении системного оптимума при сохранении автономности пользователей.

\textbf{Конкретная инстанция предметной области.} В данной работе общий класс задач рассматривается на материале формирования программы научно-практических IT-конференций. Выбор инстанции обусловлен явной выраженностью характерных черт класса: значительной неоднородностью предпочтений участников, существенным разбросом вместимости параллельных сессий, отсутствием полной информации о пользователях, а также наличием доступа к экспериментальной площадке. Эмпирическая часть работы выполняется на данных конференций Mobius и Heisenbug компании JUG.

\section{Сущность решаемой проблемы}

\textbf{Деградация системы при отсутствии управления.} В отсутствие координирующего вмешательства распределение пользователей по ресурсам формируется как агрегат независимых индивидуальных выборов. Поскольку индивидуальные предпочтения неоднородны, но имеют общие пики (наиболее популярные позиции), результирующее распределение характеризуется одновременной перегрузкой части ресурсов и хронической недозагрузкой другой части. В терминах теории игр такое состояние — равновесие эгоистичных агентов, отстоящее от системного оптимума на величину, измеряемую показателем «цена анархии» (Price of Anarchy).

\textbf{Структура потерь.} Потери при подобном состоянии затрагивают всех участников взаимодействия. Конкретный их характер варьируется в зависимости от природы домена; ниже описаны категории, имеющие универсальный характер для всего рассматриваемого класса задач.

Пользователи, обращающиеся к перегруженным ресурсам, сталкиваются с отказами в обслуживании, увеличением времени ожидания, снижением качества обслуживания, отсутствием равноценных альтернатив в момент возникновения конфликта. Конкретное стоимостное выражение этих издержек определяется доменом: прямые денежные затраты, не компенсируемые получаемой полезностью; потеря времени; упущенные возможности; ухудшение качества предоставляемой услуги.

Управляющая сторона несёт потери, связанные с неэффективным использованием инфраструктурного капитала: затраты на обеспечение недозагруженных ресурсов не компенсируются объёмом фактического использования. Снижение интегрального качества обслуживания одновременно отражается на показателях лояльности пользователей и продолжительности их участия в системе. При наличии у управляющей стороны коммерческих интересов, опосредованных распределением пользователей — например, интересов внешних партнёров, обусловленных охватом отдельных сегментов системы, — неравномерность распределения дополнительно снижает соответствующую коммерческую ценность.

В доменах, где роль поставщика ресурсов выделена как самостоятельная, поставщики недозагруженных ресурсов сталкиваются с обесцениванием своих ресурсов как канала распространения предоставляемого через них содержания. Это снижает мотивацию поставщиков к продолжению участия в системе и в долгосрочной перспективе ведёт к снижению разнообразия и качества предлагаемых ресурсов.

\textbf{Несостоятельность прямого применения стандартных подходов.} Прямое применение классической рекомендательной системы не решает поставленной задачи. Классическая рекомендательная система оптимизирует индивидуальную полезность пользователя на основании его предпочтений, не учитывая ограничений вместимости и взаимосвязанности рекомендаций. При одновременной выдаче рекомендаций множеству пользователей подобная система воспроизводит одинаково популярные позиции в верхних строках выдачи, тем самым усиливая, а не ослабляя нагрузку на наиболее востребованные ресурсы.

Альтернативная стратегия централизованного принудительного назначения также неприменима. С одной стороны, она входит в противоречие с автономностью пользователей и обесценивает само участие, основанное на свободе выбора. С другой стороны, формальное назначение не гарантирует фактического следования предписанию: реальное поведение пользователей сохраняет вероятностный характер вне зависимости от формы предъявления рекомендации.

\textbf{Природа сложности задачи.} Сложность задачи определяется одновременным выполнением трёх требований, каждое из которых разрешимо в отдельности, а в совокупности образующих противоречивую систему:

\begin{enumerate}
  \item Соблюдение индивидуальной релевантности: каждая выдаваемая рекомендация должна оставаться приемлемой для пользователя с точки зрения его предпочтений; в противном случае пользователь игнорирует подсказку.
  \item Соблюдение ограничений вместимости: суммарный ожидаемый поток пользователей на каждый ресурс не должен превышать его вместимости.
  \item Согласованность совокупной выдачи: множество персональных рекомендаций образует когерентное распределение; рекомендация одного пользователя учитывает рекомендации, выдаваемые другим.
\end{enumerate}

Третье требование принципиально отличает задачу от классической постановки рекомендательных систем, в которой пользователи рассматриваются как независимые единицы анализа. В рассматриваемом классе задач пользовательские рекомендации становятся взаимозависимыми, что переводит задачу из множества независимых ранжирований в задачу совместной оптимизации с ограничениями. Объём пространства решений и вычислительная сложность возрастают при этом нелинейно с ростом числа пользователей и ресурсов.

Указанная взаимосвязанность представляет основной источник сложности задачи и обусловливает необходимость отдельного методологического подхода, не сводимого к прямому применению существующих методов классических рекомендательных систем.

\section{Обзор существующих подходов и решений}

Описанный в подразделах 1.1–1.2 класс задач допускает постановку в нескольких самостоятельных формальных рамках, каждая из которых предлагает собственный понятийный аппарат, классы алгоритмов и методические инструменты. Ниже последовательно рассматриваются пять рамок: исследование операций, теория игр, рекомендательные системы с ограничениями, обучение с подкреплением и стохастическая оптимизация. Рамки не являются взаимоисключающими: в практической работе они пересекаются и комбинируются. Конкретный методологический синтез, выбранный в работе, представлен в главе 2.

\subsection{Постановка в рамках исследования операций}

Исследование операций (Operations Research) рассматривает задачу как комбинаторную задачу распределения дискретных пользователей по ограниченным ресурсам с целью максимизации суммарной полезности. Каноническая формулировка — обобщённая задача о назначениях (Generalized Assignment Problem, GAP), являющаяся классическим объектом изучения комбинаторной оптимизации.

Формальная постановка имеет следующий вид. Для множества пользователей $I = \{1, \ldots, n\}$, множества ресурсов $J = \{1, \ldots, m\}$, функции полезности $u_{ij}$ и вместимостей $c_j$ требуется найти бинарные индикаторы назначения $x_{ij} \in \{0, 1\}$, максимизирующие суммарную полезность $\sum_{ij} u_{ij} x_{ij}$ при условии однократного назначения каждого пользователя ($\sum_j x_{ij} = 1$ для всех $i$) и соблюдения вместимостей ($\sum_i x_{ij} \leq c_j$ для всех $j$). В случае, когда пользователю необходимо последовательно использовать несколько ресурсов в различные временные интервалы, постановка расширяется: ограничение однократного назначения заменяется условием, согласно которому в каждый момент времени за пользователем закреплено не более одного ресурса. Подобная расширенная формулировка относится к классу многопериодных задач о назначениях и пересекается с задачами составления расписаний (timetabling), для которых разработан самостоятельный методический аппарат.

С алгоритмической точки зрения GAP относится к классу NP-полных задач, что в общем случае исключает существование точного решения в полиномиальное время. Практически применяемые методы включают методы ветвей и границ для точного решения экземпляров умеренной размерности, лагранжеву релаксацию для построения оценок и приближённых решений [Fisher, Jaikumar, Van Wassenhove, 1986], а также эвристики на основе жадных алгоритмов и локального поиска. Систематический обзор алгоритмов представлен в работе Cattrysse и Van Wassenhove [1992]. Каноническим источником, охватывающим как задачи о рюкзаке, так и задачи о назначениях, является монография Martello и Toth [1990].

В прикладном плане аппарат исследования операций применяется в задачах распределения трудовых ресурсов (назначения работников на задачи с учётом нагрузки), в задачах распределения учащихся по образовательным группам, в задачах планирования производственных мощностей. В каждом из этих приложений решение задачи трактуется как итоговое распределение, подлежащее исполнению.

Рамка исследования операций даёт строгий формализм и развитый алгоритмический арсенал. Применительно к рассматриваемой задаче ограничивается допущением о детерминированном исполнении решения: в постановке GAP $x_{ij} = 1$ означает фактическое назначение, тогда как в работе решение управляющей стороны имеет статус рекомендации с вероятностным откликом. Прямое применение GAP даёт решения, оптимальные в детерминированной модели, но не учитывающие отклика пользователей.

\subsection{Постановка в рамках теории игр}

Теория игр рассматривает задачу как ситуацию совместного принятия решений автономными агентами, каждый из которых преследует собственную выгоду, а коллективное распределение возникает как результат отдельных индивидуальных выборов. Естественной формальной рамкой для рассматриваемого класса задач является игра с перегрузкой (congestion game) — класс игр, в которых полезность игрока на ресурсе зависит от числа других игроков, выбравших тот же ресурс.

Каноническая формулировка дана Розенталем [Rosenthal, 1973]. Задаётся множество игроков $N$, множество ресурсов $E$ и функция издержек $c_e: \mathbb{N} \to \mathbb{R}$ для каждого ресурса $e$, зависящая от числа игроков, выбравших этот ресурс. Стратегия игрока $i$ — подмножество ресурсов $s_i \subseteq E$. Издержки игрока определяются как $\sum_{e \in s_i} c_e(n_e(s))$, где $n_e(s)$ — число игроков, использующих ресурс $e$ в совместной стратегии $s$. Розенталем установлено, что любая такая игра обладает равновесием Нэша в чистых стратегиях и относится к классу точных потенциальных игр: существует единая функция потенциала $\Phi(s) = \sum_e \sum_{k=1}^{n_e(s)} c_e(k)$, монотонно изменяющаяся при любом одностороннем улучшении стратегии игрока. Концепция потенциальных игр впоследствии обобщена в работе Мондерера и Шэпли [Monderer, Shapley, 1996].

В задачах с большим числом однотипных пользователей, в которых индивидуальный вклад одного пользователя пренебрежимо мал по отношению к общему потоку, применяется непрерывная постановка, восходящая к работе Уордропа [Wardrop, 1952] о моделировании транспортных потоков. Согласно первому принципу Уордропа, в состоянии равновесия все используемые маршруты между парой источник–назначение имеют одинаковую стоимость, не превышающую стоимости любого неиспользуемого маршрута. Второй принцип Уордропа описывает системно-оптимальное распределение, минимизирующее суммарные издержки. Разрыв между равновесием эгоистичных пользователей и системным оптимумом является центральной характеристикой задачи.

Количественная мера данного разрыва — показатель «цена анархии» (Price of Anarchy), введённый в работе Кутсупиаса и Пападимитриу [Koutsoupias, Papadimitriou, 1999] и определяемый как отношение стоимости наихудшего равновесия к стоимости системного оптимума. Для задач эгоистичной маршрутизации с линейными функциями задержек верхняя граница цены анархии составляет 4/3 [Roughgarden, Tardos, 2002]. Иллюстрацией непредсказуемости системного поведения при индивидуальном выборе служит парадокс Брэсса [Braess, 1968], показывающий, что добавление дополнительного маршрута в сеть может ухудшить равновесное распределение и увеличить издержки всех пользователей. Систематическое изложение результатов теории эгоистичной маршрутизации представлено в монографии Раугардена [Roughgarden, 2005].

В прикладном плане теоретико-игровая постановка наиболее активно используется в задачах маршрутизации транспортных потоков, в которых навигационные сервисы реализуют рассматриваемый сценарий непосредственно: каждый пользователь стремится минимизировать собственное время в пути, а заторы возникают как естественное следствие отсутствия координации между пользователями.

Теоретико-игровая рамка естественно соответствует природе задачи (автономные пользователи, децентрализованные решения, агрегатное распределение); аппарат равновесий и цены анархии даёт теоретическую базу для анализа устойчивых состояний системы. Применительно к управляющей задаче эта рамка остаётся описательной — она отвечает на вопрос, какое распределение установится, но не даёт алгоритма построения смещающего воздействия. Переход к проектированию управляющих воздействий относится к смежным направлениям (mechanism design, information design), не предоставляющим прямых вычислительных алгоритмов для рекомендательных задач.

\subsection{Постановка в рамках рекомендательных систем с ограничениями}

Подход рекомендательных систем рассматривает задачу как формирование индивидуального ранжированного списка для каждого пользователя на основании предсказанной полезности, при этом классическая постановка дополняется системными ограничениями на агрегированное распределение выдач. В отличие от классических задач ранжирования, оптимизирующих исключительно индивидуальную релевантность, постановки с ограничениями вводят дополнительные критерии, относящиеся к свойствам совокупной выдачи: соблюдение распределительных пропорций, сохранение справедливости по отношению к поставщикам ресурсов, соблюдение системных ограничений вместимости.

Ранжирование с ограничениями (constrained ranking) — формальная постановка, систематически разработанная в работах Сингха и Йоахимса [Singh, Joachims, 2018; 2019]. Политика ранжирования $\pi$, отображающая пользователя в распределение по позициям выдачи, оптимизируется с целью максимизации ожидаемой пользовательской полезности при дополнительных линейных ограничениях на агрегированную экспозицию ресурсов:

\[
\max_\pi \mathbb{E}_{u} \left[ \sum_i \pi_i(u) r_{ui} \right] \quad \text{при} \quad \sum_{u} \pi_i(u) \in [c_i^{\min}, c_i^{\max}] \ \forall i,
\]

где $r_{ui}$ — предсказанная релевантность ресурса $i$ для пользователя $u$, а $c_i^{\min}, c_i^{\max}$ — нижние и верхние границы совокупной экспозиции ресурса $i$. Авторами показано, что задача допускает эффективное решение методами линейного программирования и применима как для оптимизации справедливости экспозиции, так и для соблюдения системных ограничений на распределение.

Калиброванные рекомендации (calibrated recommendations), предложенные Стеком [Steck, 2018], представляют альтернативную перспективу: для каждого пользователя распределение рекомендованных ресурсов по категориям должно соответствовать его историческому профилю. Постановка опирается на минимизацию расхождения Кульбака–Лейблера между распределением рекомендаций и пользовательским профилем как дополнительного слагаемого в целевой функции. Хотя задача калибровки сформулирована на уровне индивидуальной выдачи, аппарат расхождения распределений применяется для измерения системных свойств совокупной выдачи, в том числе равномерности заполнения ресурсов.

Концепция справедливости со стороны поставщиков (provider-side fairness) разработана в направлении исследований многосторонних рекомендательных систем [Abdollahpouri, Burke, Mobasher, 2020]. В её рамках в качестве заинтересованных сторон признаются не только пользователи, но и поставщики ресурсов, ожидающие равномерного или пропорционального распределения внимания. Классические работы данного направления включают Биегу, Гуммади и Вайкума [Biega, Gummadi, Weikum, 2018] о равенстве внимания (equity of attention), Целике и соавторов [Zehlike et al., 2017] о справедливом ранжировании FA*IR, а также позднейшие обзоры [Naghiaei, Rahmani, Deldjoo, 2022], систематизирующие подходы к учёту провайдерских интересов в системах ранжирования.

В прикладном плане рекомендательные системы с ограничениями применяются в стриминговых сервисах для управления экспозицией содержания (Netflix, Spotify), в торговых площадках для распределения внимания между продавцами, в задачах формирования ленты социальных сетей для управления разнообразием выдачи.

Рекомендательная рамка даёт развитый аппарат операциональных метрик (NDCG, MRR, ILD, метрики калибровки и справедливости экспозиции) и алгоритмы переранжирования, применимые поверх базовой модели релевантности. Главное ограничение для рассматриваемой задачи — формулировка ограничений как математического ожидания по распределению пользователей (выравнивание «в среднем» по большому потоку выдач), тогда как в работе требуется соблюдение вместимости в каждой конкретной коллективной выдаче. Адаптация классических методов под одноразовый совместный режим требует существенных модификаций.

\subsection{Постановка в рамках обучения с подкреплением}

Подход обучения с подкреплением (reinforcement learning, RL) рассматривает задачу как последовательный процесс взаимодействия между управляющей политикой и пользователями, в котором политика формирует рекомендательные подсказки, наблюдает отклик и адаптирует своё поведение на основании накопленного опыта взаимодействия. В отличие от подходов, оптимизирующих заранее зафиксированную целевую функцию на статической модели предпочтений, рамка обучения с подкреплением позволяет обучать политику непосредственно из взаимодействий с пользователями или их моделью, учитывая отложенные эффекты решений и динамику пользовательского поведения.

Каноническая формализация — марковский процесс принятия решений (Markov Decision Process, MDP). Управляющий агент в каждом шаге наблюдает состояние $s_t$ (контекст пользователя, история взаимодействия, состояние системы), выбирает действие $a_t$ (рекомендацию или ранжированный список), получает мгновенную награду $r_t$ (показатель полезности — клик, удовлетворение, продолжительность взаимодействия) и переходит в новое состояние $s_{t+1}$ согласно вероятности перехода $P(s_{t+1} | s_t, a_t)$. Целью является нахождение политики $\pi$, максимизирующей дисконтированную суммарную награду $\mathbb{E}_\pi \left[ \sum_{t=0}^\infty \gamma^t r_t \right]$.

Постановка задачи с ограничениями — марковский процесс с ограничениями (Constrained MDP), систематически разработанная в монографии Альтмана [Altman, 1999]. К стандартной MDP добавляется вектор стоимостных функций $c_1, \ldots, c_k$; ограничения принимают вид $\mathbb{E}_\pi \left[ \sum_t \gamma^t c_{i,t} \right] \leq d_i$. Альтманом доказано существование оптимальной стационарной стохастической политики и сводимость задачи к линейному программированию на пространстве распределений по парам «состояние–действие». В терминах рассматриваемого класса задач функции $c_i$ соответствуют ограничениям вместимости ресурсов, а пороговые значения $d_i$ — допустимой загрузке.

В качестве промежуточной рамки между классическим обучением с подкреплением и комбинаторной оптимизацией используются многорукие бандиты с ограничениями ресурсов (bandits with knapsacks), формализованные в работах Баданидиюру и соавторов [Badanidiyuru et al., 2013]. В этой постановке каждый раунд рассматривается как одно взаимодействие с пользователем без учёта долгосрочных переходов состояний, а глобальные ограничения формулируются как ограничения на суммарное потребление ресурсов. Линейная контекстная версия задачи [Agrawal, Devanur, 2016] обеспечивает алгоритмы с гарантированной верхней оценкой сожаления порядка $\tilde{O}(\sqrt{T})$.

Применение обучения с подкреплением в рекомендательных системах развивается с середины 2000-х; современный этап представлен работами о глубоком обучении с подкреплением для рекомендаций [Zheng et al., 2018], производственным применением метода REINFORCE с коррекцией смещения вне политики на крупной видеоплатформе [Chen et al., 2019] и систематическими обзорами направления [Afsar et al., 2022]. Постановки с ограничениями получили развитие в работе [Zheng et al., 2022] и в технических отчётах индустриальных лабораторий [Netflix Tech Blog, 2022], демонстрирующих жизнеспособность подхода в промышленных условиях. Большинство открытой литературы опирается на массивные объёмы накопленных данных взаимодействий, обеспечивающие надёжное оффлайн-обучение и валидацию методами оценки вне политики.

Рамка обучения с подкреплением учитывает динамику взаимодействия и отложенные эффекты решений; марковский процесс с ограничениями обеспечивает формальную базу учёта системных ограничений. Для рассматриваемой задачи характерны два ограничения. Первое — высокие требования к объёму данных: производственные применения опираются на миллионы взаимодействий, что недоступно при малом каталоге; следствие — необходимость симулятора пользовательского поведения и перенос части валидационной нагрузки в плоскость качества симулятора. Второе — формулировка ограничений как математического ожидания по политике, означающая, как и в рекомендательной рамке, выравнивание «в среднем», а не соблюдение требований в каждой коллективной выдаче.

\subsection{Постановка в рамках стохастической оптимизации}

Подход стохастической оптимизации рассматривает задачу с явным учётом случайности пользовательского поведения. Каноническими формулировками являются двухэтапное стохастическое программирование [Birge, Louveaux, 2011], в котором решение принимается в условиях неопределённости и корректируется после её раскрытия; вероятностные ограничения [Charnes, Cooper, 1959], формулируемые как требования на вероятность нарушения; робастная оптимизация [Ben-Tal, El Ghaoui, Nemirovski, 2009], дающая гарантии работоспособности по наихудшему сценарию. Прикладным образцом служит область управления доходностью (revenue management) [Talluri, van Ryzin, 2004], структурно близкая к рассматриваемой задаче.

В настоящей работе стохастическая рамка не выступает основным методологическим аппаратом; элементы вероятностного описания используются в формализации задачи (раздел 2.1), но операциональным механизмом построения политики выбран марковский процесс принятия решений с ограничениями. Высокая вычислительная сложность точных стохастических постановок и необходимость в данных для оценивания распределения спроса ограничивают их прямое применение в условиях малого каталога ресурсов и отсутствия исторических данных взаимодействий.

\section{Обзор технологий и методов}

Если в подразделе 1.3 систематизированы концептуальные рамки постановки задачи, то здесь приводится обзор конкретных алгоритмов, моделей и подходов, разработанных в литературе и индустрии. Цель — определить инструментарий, доступный для непосредственного использования или адаптации, и оценить его применимость с учётом специфики: малый каталог ресурсов, отсутствие исторических данных взаимодействий, требование соблюдения ограничений в одноразовой коллективной выдаче. Обзор проводится по четырём направлениям: алгоритмы переранжирования с системными ограничениями (1.4.1), методы обучения с подкреплением для рекомендательных задач (1.4.2), параметрические симуляционные среды (1.4.3), подходы на основе больших языковых моделей и автономных агентов (1.4.4).

\subsection{Алгоритмы переранжирования с системными ограничениями}

Переранжирование — наиболее доступный класс инструментов учёта системных ограничений в рекомендательной выдаче: к базовому ранжированию по индивидуальной релевантности применяется преобразование, обеспечивающее заданные свойства совокупного результата (разнообразие, справедливость, баланс категорий, ограничения вместимости). Базовая модель не переобучается, что упрощает внедрение в существующие системы.

Канонические алгоритмы класса включают Maximal Marginal Relevance [Carbonell, Goldstein, 1998] — жадный метод балансирования релевантности и разнообразия через параметр $\lambda$; детерминантные точечные процессы [Kulesza, Taskar, 2012], выбирающие подмножество с вероятностью, пропорциональной определителю подматрицы ядра, что штрафует похожие элементы; алгоритм FA*IR [Zehlike et al., 2017] для справедливого формирования top-k выдачи с гарантиями представления защищённой группы; алгоритм калиброванных рекомендаций [Steck, 2018], минимизирующий расхождение Кульбака–Лейблера между распределением выдачи по категориям и заданным эталонным распределением. Последний наиболее близок к задачам соблюдения распределительных ограничений на уровне коллективной выдачи: при надлежащем определении эталона аппарат расхождения распределений непосредственно применим к контролю заполненности ресурсов.

Общее ограничение алгоритмов переранжирования применительно к рассматриваемой задаче — их статический характер. Алгоритм оптимизирует выдачу при заданных предположениях о пользовательском поведении, но не учитывает вероятностного отклика и взаимной зависимости рекомендаций между пользователями; распространение на коллективную оптимизацию требует существенных обобщений, выходящих за рамки оригинальных постановок.

\subsection{Методы обучения с подкреплением для рекомендательных задач}

Методы обучения с подкреплением, применяемые в рекомендательных системах, делятся на три семейства: основанные на функции ценности (Q-обучение [Watkins, Dayan, 1992] и его глубокая версия DQN [Mnih et al., 2015], применённая в DRN [Zheng et al., 2018] для рекомендации новостей); основанные на политике (REINFORCE [Williams, 1992]); гибридные методы «актор-критик» (A2C/A3C [Mnih et al., 2016] и PPO [Schulman et al., 2017]). Семейство «актор-критик» естественно интегрируется с задачами с ограничениями через двойственную оптимизацию множителей Лагранжа [Zheng et al., 2022].

Производственные применения опираются на накопленные журналы взаимодействий и требуют методов компенсации смещения между обучающей и целевой политиками — выборки по значимости (importance sampling) и её устойчивых вариантов, а также методов дважды робастной оценки. Работа Чена и соавторов [Chen et al., 2019] предложила Top-K коррекцию вне политики для метода REINFORCE, обеспечив развёртывание обучаемой политики в крупной производственной системе.

В рассматриваемой задаче прямое применение обучения с подкреплением осложнено отсутствием накопленных журналов взаимодействий: малый объём исторических данных делает методы оценки вне политики ненадёжными. Альтернативой служит обучение и оценка политик в симуляционной среде (подраздел 1.4.3).

\subsection{Параметрические симуляционные среды для рекомендательных систем}

При недоступности достаточных объёмов накопленных данных взаимодействий и невозможности масштабных онлайн-экспериментов основным альтернативным источником данных служит симуляционная среда, моделирующая поведение пользователей в ответ на рекомендательные политики. Параметрический класс симуляторов задаёт поведение через явно специфицированные вероятностные функции выбора; альтернативный класс на основе больших языковых моделей рассматривается в подразделе 1.4.4.

Каноническими представителями класса являются RecSim [Ie et al., 2019] (Google; модель пользователя с латентным состоянием, моделями выбора и интеграцией с библиотеками обучения с подкреплением); Sim4Rec [Sber AI Lab, 2022] (открытая разработка ПАО «Сбер», ориентированная на автономные эксперименты с поддержкой синтетической генерации пользователей и моделей выбора, имеющая документацию на русском языке); KuaiSim [Zhao et al., 2023] (среда исследовательской группы Kuaishou с поведенческой моделью, обученной на журналах производственной системы) с сопровождающим набором полностью наблюдаемых данных KuaiRec [Gao et al., 2022].

Параметрические симуляционные среды применимы к рассматриваемой задаче, однако сопровождаются известными методологическими оговорками: разрыв между симулируемым и реальным поведением пользователей (sim-to-real gap) ставит вопрос внешней валидности; сам симулятор содержит содержательные предположения о поведении, требующие явного обоснования; при использовании среды как основной площадки сравнения методов необходима независимая проверка устойчивости результатов к выбору параметров.

\subsection{Подходы на основе больших языковых моделей и автономных агентов}

В 2023–2026 годах сформировался самостоятельный класс подходов, основанных на использовании больших языковых моделей (large language models, LLM) и автономных агентов на их базе. В отличие от параметрических методов, агентный подход возлагает функцию принятия решений на экземпляр LLM, получающей описание персоны пользователя в естественном языке. Большая языковая модель может выступать также в роли ранжирующего компонента, генератора синтетических данных, инструмента оценки качества или элемента контура обучения.

В роли симуляторов пользователей идея восходит к работе о генеративных агентах [Park et al., 2023] и развивается в серии работ для рекомендательного домена — Agent4Rec [Zhang et al., 2024], AgentCF, RecAgent — а также в крупномасштабных средах OASIS [Yang et al., 2024] и AgentSociety [Piao et al., 2025], моделирующих миллионы агентов в социальных и городских контекстах.

В роли ранжирующего компонента современные LLM способны выполнять ранжирование без дополнительного обучения (zero-shot) [Hou et al., 2024], интегрируясь поверх существующих систем поиска по эмбеддингам. Семейство методов LLMRec, P5, GPT4Rec и ChatRec, а также свежие подходы BIGRec [Bao et al., 2024] и LLaRA [Liao et al., 2024] исследуют разные способы использования LLM в задачах рекомендации — от прямой генерации списков до гибридных архитектур, объединяющих эмбеддинги традиционных рекомендательных моделей с языковыми моделями. Агентные рекомендательные системы развивают подход через паттерны ReAct [Yao et al., 2022] и Reflexion.

В задачах генерации синтетических данных и оценки качества LLM применяются для создания пользовательских профилей и сценариев поведения, что особенно ценно для задач с малым каталогом ресурсов. Парадигма LLM-as-judge [Zheng et al., 2023] обеспечивает ускоренную автономную оценку рекомендательных политик без дорогостоящих пользовательских исследований.

Гибридные схемы комбинируют преимущества обоих классов: однократная LLM-генерация персон с последующим использованием в параметрическом симуляторе существенно снижает совокупную вычислительную стоимость при сохранении содержательного разнообразия популяции. Двухслойная валидация — обучение в параметрическом симуляторе и независимая проверка в LLM-агентной среде — согласует требования эффективности и достоверности результатов. Главные ограничения подхода — вычислительная стоимость прямого использования LLM на каждом шаге и вопрос соответствия поведения LLM-агентов поведению реальных пользователей, активно обсуждаемый в свежей литературе.

\section{Анализ ограничений существующих решений}

Систематизация рассмотренных в подразделах 1.3–1.4 концептуальных рамок, методов и инструментов позволяет выделить ключевые ограничения существующих подходов. Анализ проводится по нескольким независимым осям, каждая из которых отражает важное свойство задачи, недостаточно покрываемое существующей литературой.

\textbf{Ось 1. Совместная оптимизация рекомендаций между пользователями.} Большинство рамок рассматривают задачу на уровне индивидуального пользователя; постановка, в которой множество пользователей одновременно требуют согласованной выдачи, представлена лишь в исследовании операций и в специализированных постановках управления доходностью, причём первое не учитывает вероятностного отклика.

\textbf{Ось 2. Жёсткие ограничения вместимости в одноразовой выдаче.} Подходы рекомендательных систем с ограничениями и постановки Constrained MDP формулируют ограничения как математическое ожидание по распределению пользователей или по политике, что обеспечивает выравнивание «в среднем» по большому потоку выдач, но не соблюдение в каждой конкретной коллективной выдаче, требуемое в рассматриваемой постановке.

\textbf{Ось 3. Ограниченность исторических данных.} Производственные применения опираются на миллионы накопленных взаимодействий; в задачах с малым каталогом и ограниченным числом пользователей подобный объём данных принципиально недоступен, что исключает прямое применение методов обучения с подкреплением и методов оценки вне политики.

\textbf{Ось 4. Природа управляющего рычага.} Классические постановки покрывают полярные случаи — либо детерминированное исполнение, либо отсутствие управляющего воздействия. Промежуточный случай мягкой рекомендации, которой пользователь может не следовать, формализован слабее; смежные направления (mechanism design, information design) развиты для теоретического анализа и не предоставляют прямых вычислительных алгоритмов.

\textbf{Ось 5. Симуляционная валидация при ограниченных данных.} Параметрические симуляторы (RecSim, Sim4Rec, KuaiSim) ориентированы на сценарии с большими каталогами и параметрическими моделями предпочтений, что несёт риск упрощения при малом каталоге. LLM-агентные симуляторы (подраздел 1.4.4) дают содержательную глубину, но сопровождаются высокой вычислительной стоимостью и открытым вопросом соответствия реальному поведению пользователей.

\textbf{Идентифицированный пробел.} Объединение перечисленных осей формирует область, систематически не покрытую существующими решениями: задача формирования согласованных рекомендательных выдач для множества пользователей одновременно, с жёсткими ограничениями вместимости в каждой одноразовой выдаче, при ограниченных исторических данных, в условиях мягкого характера управляющего рычага и при необходимости опоры на симуляционную валидацию. Привлечение средств LLM-агентного аппарата (подраздел 1.4.4) существенно расширяет возможности работы с осями 3 и 5 (ограниченность данных и симуляционная валидация), однако само по себе не закрывает оси 1 и 2 (совместная оптимизация и жёсткие ограничения вместимости в одноразовой выдаче). Ни одна из рассмотренных концептуальных рамок и ни один из существующих программных инструментов — включая инструменты, основанные на больших языковых моделях, — не предоставляет готового решения для указанного пересечения требований; работы, систематически адресующие все пять осей одновременно, в открытой литературе на момент проведения исследования не обнаружены. Этот пробел определяет содержательное пространство для разработки методологического и алгоритмического подхода в работе; конкретное наполнение пространства представлено в главе 2.

\section{Выводы по главе}

Рассматриваемый класс задач — распределение автономных пользователей по ресурсам ограниченной вместимости через рекомендательную подсказку, не носящую принудительного характера, — возникает в широком спектре прикладных областей и образует устойчивую предметную постановку. Конкретной инстанцией работы выбрано формирование программы научно-практических IT-конференций.

Сложность задачи определяется одновременным выполнением трёх требований: соблюдения индивидуальной релевантности рекомендаций, соблюдения ограничений вместимости ресурсов и согласованности совокупной выдачи между пользователями. Прямое применение классических рекомендательных систем эту задачу не закрывает, поскольку независимая оптимизация индивидуальной полезности усиливает, а не ослабляет нагрузку на наиболее востребованные ресурсы.

Среди существующих концептуальных рамок ни одна не покрывает совокупности характерных для задачи требований. Аппарат исследования операций ограничен предположением о детерминированном исполнении решения; классическая теория игр носит описательный, а не предписывающий характер; постановки рекомендательных систем с ограничениями формулируют ограничения как амортизированные ожидания, а не как требования к одной выдаче; методы обучения с подкреплением требуют объёмов данных, недоступных в условиях малого каталога; стохастическая оптимизация не предоставляет операциональной схемы построения управляющей политики. Аналогичный вывод получен для рассмотренных алгоритмов и симуляционных сред — ни один из готовых инструментов не закрывает специфики задачи без существенной адаптации.

Идентифицированный пробел, описанный в подразделе 1.5, определяет содержательное пространство работы. В главе 2 формализуется постановка задачи и описывается предлагаемый методологический подход, опирающийся на синтез аппарата марковского процесса принятия решений с ограничениями и теории игр с перегрузкой.
